\section{Introduction}
Clustering is a foundational problem in machine learning and data analysis, aiming to partition a dataset into subsets (clusters) such that data points within a cluster are similar according to a given objective. Classical clustering algorithms, such as $k$-means, $k$-median, and $k$-center, have found widespread use in applications ranging from recommendation systems and market segmentation to medical diagnostics and image analysis. However, these traditional objectives focus solely on optimizing notions of proximity or compactness and often overlook issues of \emph{fairness}, particularly in how individuals from different demographic or social groups are treated.

In recent years, there has been a growing awareness that standard clustering procedures may exacerbate bias or reinforce structural disparities, such as forming clusters that disproportionately represent one demographic group while marginalizing others. This has motivated the development of \emph{fair clustering} algorithms, which aim to incorporate fairness constraints alongside standard clustering objectives. Fairness in clustering can take many forms, including ensuring balanced group representation within clusters, minimizing worst-case costs across demographic groups (social fairness), or guaranteeing that similar individuals are treated similarly (individual fairness). These models bring together algorithmic design with ethical and societal considerations, making the study of fair clustering both technically rich and socially important.

In this survey, we provide an overview of algorithmic results in fair clustering, with an emphasis on theoretical models, approximation guarantees, and algorithmic techniques. Our focus is primarily on \emph{centroid-based clustering} objectives, such as fair variants of $k$-center, $k$-median, and $k$-means, which have been the most extensively studied from an algorithmic perspective. We briefly mention other clustering frameworks where fairness has also been explored, including \emph{spectral clustering}~\citep{DBLP:conf/icml/KleindessnerSAM19,DBLP:conf/aistats/WangLDB23,tonin2025accelerating} and \emph{hierarchical clustering}~\citep{DBLP:conf/nips/AhmadianEK0MMPV20,DBLP:conf/afci/ChhabraSM21,DBLP:conf/icml/KnittelSDH23,DBLP:conf/nips/KnittelSDH23}, though these are not the main focus of this survey.

It is worth emphasizing that this is by no means a comprehensive survey. The body of work in fair clustering is vast and growing, including numerous models, fairness criteria, algorithmic strategies, and empirical evaluations. Our goal is to present a curated overview of core algorithmic contributions that have shaped the field, particularly in terms of theoretical guarantees and complexity results.
Moreover, we are not covering the empirical aspects of fair clustering in detail, such as the development of real-world benchmarks, practical deployment challenges, and critical evaluations of fairness definitions. These components are not only important but also essential for advancing the overall impact of fair clustering within the broader fields of algorithmic fairness and machine learning. Developing representative datasets and benchmarks is crucial for translating theoretical insights into effective practice, and deployment-focused studies often reveal practical trade-offs that are not captured by algorithmic models alone. For readers interested in these empirical and application-oriented directions, we point to the growing body of work that investigates the effectiveness of fair clustering on real datasets, examines benchmark design, and assesses the interplay between fairness, accuracy, and interpretability in real-world settings~\citep{chhabra2021overview,abbasi2021fair}. Well-known real-world examples, such as the ProPublica analysis of recidivism risk scores in classification~\citep{angwin2022machine}, have been instrumental in highlighting both the challenges and societal impact of fairness in algorithmic decision making. Similarly, compelling case studies and practical applications in fair clustering have the potential to capture attention and clearly demonstrate the real-world value and necessity of fairness in this area.

\paragraph{Roadmap of the survey.}
Fair clustering methods can be broadly categorized according to the type of fairness they seek to enforce, reflecting different philosophical and practical interpretations of what it means for a clustering to be \emph{``fair''}~\citep{dwork2012fairness,hardt2016equality,barocas2017fairness,mitchell2021algorithmic}. These categories distinguish whether fairness is defined at the level of individuals, groups, or the overall distribution of clustering outcomes, and each perspective introduces unique algorithmic challenges and modeling considerations.
